\documentclass[11pt,a4paper]{article}

% ── Codificación y fuentes ────────────────────────────────────────────────────
\usepackage[utf8]{inputenc}
\usepackage[T1]{fontenc}
\usepackage{lmodern}
\usepackage[spanish]{babel}

% ── Geometría y espaciado ─────────────────────────────────────────────────────
\usepackage[top=2.5cm, bottom=2.5cm, left=2.8cm, right=2.8cm]{geometry}
\usepackage{setspace}
\setstretch{1.15}
\usepackage{parskip}

% ── Tablas ────────────────────────────────────────────────────────────────────
\usepackage{booktabs}
\usepackage{tabularx}
\usepackage{array}
\usepackage{longtable}
\usepackage{enumitem}

% ── Colores y cajas ───────────────────────────────────────────────────────────
\usepackage[dvipsnames]{xcolor}
\usepackage{tcolorbox}
\tcbuselibrary{skins, breakable}

% ── Cabeceras y pies de página ────────────────────────────────────────────────
\usepackage{fancyhdr}
\usepackage{lastpage}
\pagestyle{fancy}
\fancyhf{}
\fancyhead[L]{\small\textbf{Informe de Evaluación} --- Física para Bioanalistas}
\fancyhead[R]{\small maria melendez 33685873}
\fancyfoot[C]{\small Página \thepage\ de \pageref{LastPage}}
\renewcommand{\headrulewidth}{0.4pt}

% ── Hiperenlaces ──────────────────────────────────────────────────────────────
\usepackage[colorlinks=true, linkcolor=NavyBlue, urlcolor=NavyBlue]{hyperref}

% ── Paleta de colores personalizados ─────────────────────────────────────────
\definecolor{desempeno_excelente}{HTML}{1A6B2F}
\definecolor{desempeno_bueno}{HTML}{2E5A9C}
\definecolor{desempeno_suficiente}{HTML}{8B6914}
\definecolor{desempeno_deficiente}{HTML}{C0392B}
\definecolor{desempeno_insuficiente}{HTML}{7B241C}
\definecolor{grisFondo}{HTML}{F5F5F5}
\definecolor{azulTitulo}{HTML}{1B2A6B}
\definecolor{verdeFortaleza}{HTML}{1A6B2F}
\definecolor{naranjaMejora}{HTML}{A04000}

% ── Estilos de cajas ─────────────────────────────────────────────────────────
\tcbset{
  cajaTitulo/.style={
    enhanced, breakable,
    colback=azulTitulo!8, colframe=azulTitulo,
    fonttitle=\bfseries\large, coltitle=white,
    attach boxed title to top left={yshift=-2mm, xshift=4mm},
    boxed title style={colback=azulTitulo},
    top=6pt, bottom=6pt, left=8pt, right=8pt,
  },
  cajaFortaleza/.style={
    enhanced, breakable,
    colback=verdeFortaleza!6, colframe=verdeFortaleza!60,
    fonttitle=\bfseries, coltitle=verdeFortaleza,
    top=4pt, bottom=4pt, left=8pt, right=8pt,
  },
  cajaMejora/.style={
    enhanced, breakable,
    colback=naranjaMejora!6, colframe=naranjaMejora!60,
    fonttitle=\bfseries, coltitle=naranjaMejora,
    top=4pt, bottom=4pt, left=8pt, right=8pt,
  },
  cajaRecomendacion/.style={
    enhanced, breakable,
    colback=grisFondo, colframe=gray!50,
    fonttitle=\bfseries, coltitle=black,
    top=4pt, bottom=4pt, left=8pt, right=8pt,
  },
}

% ─────────────────────────────────────────────────────────────────────────────
\begin{document}

% ══ PORTADA ══════════════════════════════════════════════════════════════════
\begin{center}
  {\Large\bfseries\color{azulTitulo} Universidad de Oriente/Núcleo Sucre --- Física para Bioanalistas}\\[4pt]
  {\large Informe de Evaluación Preliminar}\\[12pt]
  \rule{\linewidth}{1.2pt}\\[6pt]
  {\LARGE\bfseries maria melendez 33685873}\\[4pt]
  \rule{\linewidth}{0.6pt}
\end{center}

% ══ FICHA TÉCNICA ════════════════════════════════════════════════════════════
\vspace{4pt}
\begin{tcolorbox}[cajaTitulo, title=Datos del informe]
\begin{tabular}{@{}llll@{}}
  \textbf{Estudiante:}  & maria melendez 33685873  &
  \textbf{Fecha:}       & 2026-06-26 \\[4pt]
  \textbf{Archivo PDF:} & \texttt{maria\_melendez\_33685873.pdf}  &
  \textbf{Páginas:}     & 4 \\
\end{tabular}
\end{tcolorbox}

% ══ RESULTADO GLOBAL ═════════════════════════════════════════════════════════
\vspace{6pt}
\begin{center}
  \begin{tcolorbox}[
    enhanced,
    colback=white, colframe=desempeno_bueno,
    width=0.62\linewidth,
    boxrule=2pt,
    halign=center, valign=center,
    top=8pt, bottom=8pt,
  ]
    \centering
    {\large\bfseries Puntaje sugerido}\\[6pt]
    {\Huge\bfseries\color{desempeno_bueno} 8.0/10}\\[6pt]
    {\large Nivel de desempeño: \textbf{\color{desempeno_bueno} Bueno}}
  \end{tcolorbox}
\end{center}

% ══ RESUMEN GENERAL ══════════════════════════════════════════════════════════
\section*{Resumen general}
El estudiante demuestra una buena comprensión de los principios fundamentales de la física de fluidos aplicados al bioanálisis. La mayoría de los problemas fueron abordados con la metodología correcta y la aplicación adecuada de las fórmulas. Sin embargo, se observan errores en la interpretación de una pregunta, en la transcripción de una fórmula y en la precisión de los cálculos numéricos finales, lo que afecta el puntaje global.

% ══ FORTALEZAS Y ÁREAS DE MEJORA ════════════════════════════════════════════
\vspace{4pt}
\begin{minipage}[t]{0.48\linewidth}
  \begin{tcolorbox}[cajaFortaleza, title={Fortalezas}]
\begin{itemize}[leftmargin=1.5em, itemsep=2pt]
  \item Correcta aplicación de principios físicos fundamentales (presión hidrostática, Arquímedes, continuidad, Bernoulli, Poiseuille).
  \item Habilidad para identificar datos, convertir unidades y reordenar fórmulas correctamente en la mayoría de los casos.
  \item Presentación clara y organizada de los pasos de resolución para la mayoría de los problemas.
  \item Manejo adecuado de unidades en los cálculos.
\end{itemize}
  \end{tcolorbox}
\end{minipage}
\hfill
\begin{minipage}[t]{0.48\linewidth}
  \begin{tcolorbox}[cajaMejora, title={Areas de mejora}]
\begin{itemize}[leftmargin=1.5em, itemsep=2pt]
  \item Precisión en la interpretación de los enunciados de los problemas, especialmente en términos como 'cambio porcentual'.
  \item Rigurosidad en la transcripción de fórmulas y en la presentación de los pasos de cálculo.
  \item Cuidado en los cálculos numéricos finales, prestando especial atención a la notación científica y la posición del punto decimal.
\end{itemize}
  \end{tcolorbox}
\end{minipage}

% ══ OBSERVACIONES POR PREGUNTA ═══════════════════════════════════════════════
\section*{Observaciones por pregunta}

\begin{longtable}{>{\bfseries}p{2.8cm} p{1.6cm} p{7.0cm} p{2.8cm}}
  \toprule
  \textbf{Pregunta} & \textbf{Puntaje} & \textbf{Evaluación} & \textbf{Errores detectados} \\
  \midrule
  \endhead
  \bottomrule
  \endfoot
    \textbf{1. Presión Hidrostática y Venosa} & 2.0/2.0 & El estudiante aplica correctamente la fórmula de presión hidrostática (Pm = $\rho$gh) y realiza los cálculos de manera precisa. La identificación de datos y la conversión de unidades son correctas. La respuesta final es exacta y con unidades consistentes. & \textit{---} \\
    \midrule
    \textbf{2. Principio de Arquímedes} & 1.5/2.0 & El cálculo de la fuerza de empuje (E = W\_real - W\_aparente) y el volumen del objeto (V\_muestra = E / ($\rho$\_agua * g)) son correctos. Sin embargo, la fórmula escrita para la densidad de la muestra (P\_muestra = W\_real / (1000 kg/m$^{3}$ * 1.875 * 1000)) es incorrecta y confusa, aunque el resultado numérico final (1875 kg/m$^{3}$) es, por coincidencia, el valor correcto. Se penaliza la falta de claridad y la incorrecta transcripción de la fórmula. & \textit{Fórmula de densidad escrita incorrectamente, a pesar de que el resultado numérico final es correcto.} \\
    \midrule
    \textbf{3. Hemodinámica (Continuidad y Bernoulli)} & 2.0/2.0 & El estudiante aplica correctamente la ecuación de continuidad para calcular la velocidad en la zona estrecha (V2) y luego utiliza la ecuación de Bernoulli para determinar la presión (P2). Todos los pasos, sustituciones y cálculos son precisos. La organización es excelente. & \textit{---} \\
    \midrule
    \textbf{4. Ley de Poiseuille (Análisis Proporcional)} & 1.0/2.0 & El estudiante identifica correctamente la relación Q $\propto$ r$^{4}$ y calcula la nueva relación de flujo (Q'/Q = (0.80)$^{4}$ = 0.4096). Sin embargo, interpreta erróneamente la pregunta al dar el 'porcentaje de flujo remanente' (40.96\%) en lugar del 'porcentaje de cambio' en el flujo, que debería ser una disminución del 59.04\%. & \textit{Error conceptual en la interpretación de 'porcentaje de cambio' en el flujo.} \\
    \midrule
    \textbf{5. Flujo Viscoso y Resistencia} & 1.5/2.0 & El estudiante identifica correctamente todos los datos, convierte las unidades adecuadamente y aplica la Ley de Poiseuille (Q = ($\pi$R$^{4}$$\Delta$P) / (8$\eta$L)) para despejar la diferencia de presión ($\Delta$P). La sustitución de valores es correcta. Sin embargo, el resultado numérico final (1.989.43 Pa) presenta un error significativo en la magnitud (un factor de 100), ya que el valor correcto es aproximadamente 198943.6 Pa. & \textit{Error procedimental significativo en el cálculo numérico final (error de magnitud/punto decimal).} \\
    \midrule
\end{longtable}

% ══ ERRORES DETECTADOS ═══════════════════════════════════════════════════════
\section*{Análisis de errores}

\subsection*{Errores conceptuales}
\begin{itemize}[leftmargin=1.5em, itemsep=2pt]
  \item Confusión entre 'flujo remanente como porcentaje' y 'cambio porcentual' en la Pregunta 4.
\end{itemize}

\subsection*{Errores procedimentales}
\begin{itemize}[leftmargin=1.5em, itemsep=2pt]
  \item Transcripción incorrecta de la fórmula de densidad en la Pregunta 2, a pesar de obtener el resultado numérico correcto.
  \item Error de cálculo numérico (magnitud/punto decimal) en la respuesta final de la Pregunta 5.
\end{itemize}

% ══ RECOMENDACIONES AL ESTUDIANTE ════════════════════════════════════════════
\section*{Retroalimentación para el estudiante}
\begin{tcolorbox}[cajaRecomendacion, title={Dirigido a: maria melendez 33685873}]
Estimado estudiante, has demostrado una sólida base en los principios de la física de fluidos. Para mejorar aún más, te sugiero revisar cuidadosamente la redacción de las preguntas, especialmente términos como 'cambio porcentual', para asegurar que tu respuesta final refleje exactamente lo solicitado. Presta mucha atención a la precisión en los cálculos numéricos, especialmente al manejar notación científica y la posición del punto decimal, ya que un pequeño error puede llevar a un resultado final incorrecto. Finalmente, asegúrate de que todas las fórmulas y pasos intermedios estén correctamente escritos; esto no solo te ayuda a ti a verificar tu trabajo, sino que también demuestra una comprensión completa del procedimiento.
\end{tcolorbox}

% ══ NOTA PARA EL PROFESOR ════════════════════════════════════════════════════
\section*{Nota para el profesor}
\begin{tcolorbox}[
  enhanced, breakable,
  colback=yellow!8, colframe=yellow!60!black,
  fonttitle=\bfseries, coltitle=black,
  title={Observaciones para revisión manual},
  top=4pt, bottom=4pt, left=8pt, right=8pt,
]
En la Pregunta 2, el estudiante escribió una fórmula incorrecta para la densidad, pero el resultado numérico final es correcto. Esto podría indicar un error de transcripción o que el estudiante llegó al resultado por otro medio. Se sugiere una revisión para entender cómo llegó a ese número. En la Pregunta 4, el error es de interpretación del término 'cambio porcentual', no de la ley física subyacente, lo cual es un error común.
\end{tcolorbox}

% ══ PIE DE INFORME ════════════════════════════════════════════════════════════
\vspace{12pt}
\begin{center}
  \footnotesize\color{gray}
  Informe generado automáticamente por el sistema de evaluación con IA (gemini-2.5-flash) y soporte LaTex. \\
  Este documento es una evaluación \textbf{preliminar} para ser revisado por el profesor antes de ser entregado al 
  estudiante. \\
  Generado el 2026-06-26 \textbullet\ BIOANALISIS Sistema Académico de Evaluación.
  \end{center}

\end{document}
